% GENERATED FILE. Edit canonical lessons, then run npm run generate:books.
% canonical-source-hash: fnv1a64:51e2f9979a7b3fba
% canonical-lessons: TE-C219-ninda, TE-C219-adharam, TE-C219-rujuvu, TE-C219-parinamam, TE-C219-prerana

\chapter{Blame, Basis, Proof, Development, Motivation}
\label{ch:te-blame-basis-proof-development-motivation}
\hlchaptermodality{219}

\begin{chapteropening}
\textbf{By the end of this chapter:} \emph{I can say blame, a basis, proof, a development and motivation.}

Complete the last lesson of chapter 219: I can say blame, a basis, proof, a development and motivation.
\end{chapteropening}

\section[\texorpdfstring{\te{నింద} (ninda)}{నింద (ninda)}]{\te{నింద} (ninda) — blame}

\label{lesson:TE-C219-ninda}



\begin{quote}
Before the new one: say the Telugu for a difference, then the Telugu for effort.
\end{quote}

\subsection*{You'll want to know: \te{నింద}}
\textbf{\te{నింద}} — \emph{ninda} — \textquotedblleft{}blame\textquotedblright{}.

\begin{grammarlens}[title={What you've built}]
One more word for this chapter.
\end{grammarlens}

\subsection*{Your turn}
\begin{itemize}
\raggedright
  \item \emph{Say it:} \emph{ninda}
  \item \emph{Say it:} \emph{ninda}, once more
  \item \emph{Say it:} say \emph{ninda}
  \item \emph{Recall:} say the Telugu for a difference, then the Telugu for effort, then say \emph{ninda} again
\end{itemize}

\begin{tcolorbox}[breakable,colback=teal!4,colframe=teal!35!black,title={Before you move on}]
What does \te{నింద} mean? (\textquotedblleft{}Blame\textquotedblright{}.) Say it once more.
\end{tcolorbox}
\section[\texorpdfstring{\te{ఆధారం} (ādhāraṁ)}{ఆధారం (ādhāraṁ)}]{\te{ఆధారం} (ādhāraṁ) — a basis; evidence}

\label{lesson:TE-C219-adharam}



\begin{quote}
Before the new one: say the Telugu for effort, then the Telugu for blame.
\end{quote}

\subsection*{You'll want to know: \te{ఆధారం}}
\textbf{\te{ఆధారం}} — \emph{ādhāraṁ} — \textquotedblleft{}a basis; evidence\textquotedblright{}.

\begin{grammarlens}[title={What you've built}]
One more word for this chapter.
\end{grammarlens}

\subsection*{Your turn}
\begin{itemize}
\raggedright
  \item \emph{Say it:} \emph{ādhāraṁ}
  \item \emph{Say it:} \emph{ādhāraṁ}, once more
  \item \emph{Say it:} say \emph{ādhāraṁ}
  \item \emph{Recall:} say the Telugu for effort, then the Telugu for blame, then say \emph{ādhāraṁ} again
\end{itemize}

\begin{tcolorbox}[breakable,colback=teal!4,colframe=teal!35!black,title={Before you move on}]
What does \te{ఆధారం} mean? (\textquotedblleft{}A basis; evidence\textquotedblright{}.) Say it once more.
\end{tcolorbox}
\section[\texorpdfstring{\te{రుజువు} (rujuvu)}{రుజువు (rujuvu)}]{\te{రుజువు} (rujuvu) — proof}

\label{lesson:TE-C219-rujuvu}



\begin{quote}
Before the new one: say the Telugu for blame, then the Telugu for a basis.
\end{quote}

\subsection*{You'll want to know: \te{రుజువు}}
\textbf{\te{రుజువు}} — \emph{rujuvu} — \textquotedblleft{}proof\textquotedblright{}.

\begin{grammarlens}[title={What you've built}]
One more word for this chapter.
\end{grammarlens}

\subsection*{Your turn}
\begin{itemize}
\raggedright
  \item \emph{Say it:} \emph{rujuvu}
  \item \emph{Say it:} \emph{rujuvu}, once more
  \item \emph{Say it:} say \emph{rujuvu}
  \item \emph{Recall:} say the Telugu for blame, then the Telugu for a basis, then say \emph{rujuvu} again
\end{itemize}

\begin{tcolorbox}[breakable,colback=teal!4,colframe=teal!35!black,title={Before you move on}]
What does \te{రుజువు} mean? (\textquotedblleft{}Proof\textquotedblright{}.) Say it once more.
\end{tcolorbox}
\section[\texorpdfstring{\te{పరిణామం} (pariṇāmaṁ)}{పరిణామం (pariṇāmaṁ)}]{\te{పరిణామం} (pariṇāmaṁ) — a development, an outcome}

\label{lesson:TE-C219-parinamam}



\begin{quote}
Before the new one: say the Telugu for a basis, then the Telugu for proof.
\end{quote}

\subsection*{You'll want to know: \te{పరిణామం}}
\textbf{\te{పరిణామం}} — \emph{pariṇāmaṁ} — \textquotedblleft{}a development, an outcome\textquotedblright{}.

\begin{grammarlens}[title={What you've built}]
One more word for this chapter.
\end{grammarlens}

\subsection*{Your turn}
\begin{itemize}
\raggedright
  \item \emph{Say it:} \emph{pariṇāmaṁ}
  \item \emph{Say it:} \emph{pariṇāmaṁ}, once more
  \item \emph{Say it:} say \emph{pariṇāmaṁ}
  \item \emph{Recall:} say the Telugu for a basis, then the Telugu for proof, then say \emph{pariṇāmaṁ} again
\end{itemize}

\begin{tcolorbox}[breakable,colback=teal!4,colframe=teal!35!black,title={Before you move on}]
What does \te{పరిణామం} mean? (\textquotedblleft{}A development, an outcome\textquotedblright{}.) Say it once more.
\end{tcolorbox}
\section[\texorpdfstring{\te{ప్రేరణ} (prēraṇa)}{ప్రేరణ (prēraṇa)}]{\te{ప్రేరణ} (prēraṇa) — motivation, inspiration}

\label{lesson:TE-C219-prerana}



\begin{quote}
Before the new one: say the Telugu for proof, then the Telugu for a development.
\end{quote}

\subsection*{You'll want to know: \te{ప్రేరణ}}
\textbf{\te{ప్రేరణ}} — \emph{prēraṇa} — \textquotedblleft{}motivation, inspiration\textquotedblright{}.

\begin{grammarlens}[title={What you've built}]
One more word for this chapter.
\end{grammarlens}

\subsection*{Your turn}
\begin{itemize}
\raggedright
  \item \emph{Say it:} \emph{prēraṇa}
  \item \emph{Say it:} \emph{prēraṇa}, once more
  \item \emph{Say it:} say \emph{prēraṇa}
  \item \emph{Recall:} say the Telugu for proof, then the Telugu for a development, then say \emph{prēraṇa} again
\end{itemize}

\begin{tcolorbox}[breakable,colback=teal!4,colframe=teal!35!black,title={Before you move on}]
What does \te{ప్రేరణ} mean? (\textquotedblleft{}Motivation, inspiration\textquotedblright{}.) Say it once more.
\end{tcolorbox}
